\documentclass[11pt,reqno]{amsart}
\usepackage[localtheorems]{raeez-math-template}

% ================================================================
% Packages
% ================================================================
\usepackage{array}
\usepackage[all]{xy}

% ================================================================
% Theorem environments
% ================================================================
\newtheorem{theorem}{Theorem}[section]
\newtheorem{proposition}[theorem]{Proposition}
\newtheorem{lemma}[theorem]{Lemma}
\newtheorem{corollary}[theorem]{Corollary}
\theoremstyle{definition}
\newtheorem{definition}[theorem]{Definition}
\newtheorem{construction}[theorem]{Construction}
\newtheorem{example}[theorem]{Example}
\theoremstyle{remark}
\newtheorem{remark}[theorem]{Remark}

% ================================================================
% Macros
% ================================================================
\newcommand{\cA}{\mathcal{A}}
\newcommand{\cC}{\mathcal{C}}
\newcommand{\cD}{\mathcal{D}}
\providecommand{\cF}{\mathcal{F}}
\newcommand{\cO}{\mathcal{O}}
\newcommand{\cR}{\mathcal{R}}
\newcommand{\fg}{\mathfrak{g}}
\newcommand{\fh}{\mathfrak{h}}
\newcommand{\fn}{\mathfrak{n}}
\newcommand{\fsl}{\mathfrak{sl}}
\newcommand{\fso}{\mathfrak{so}}
\newcommand{\fsp}{\mathfrak{sp}}
\newcommand{\MC}{\mathrm{MC}}
\newcommand{\Hom}{\mathrm{Hom}}
\newcommand{\End}{\mathrm{End}}
\newcommand{\Ext}{\mathrm{Ext}}
\newcommand{\DK}{\mathrm{DK}}
\newcommand{\Rep}{\mathrm{Rep}}
\newcommand{\Omg}{\Omega}
\newcommand{\Eone}{E_1}
\newcommand{\thick}{\mathrm{thick}}
\newcommand{\ch}{\mathrm{ch}}
\newcommand{\Yh}{Y_\hbar(\fg)}
\newcommand{\Uq}{U_q(\hat{\fg})}
\newcommand{\barB}{\bar{B}}
\providecommand{\cW}{\mathcal{W}}

\DeclareMathOperator{\gr}{gr}

\numberwithin{equation}{section}

% ================================================================
\begin{document}

\title[Thin $\ell$-weight filtrations in the Drinfeld--Kohno core]
{Thin $\ell$-weight filtrations in the\\
Drinfeld--Kohno compact core}

\author{Raeez Lorgat}
\address{Department of Mathematics, Harvard University,
 Cambridge, MA 02138, USA}
\email{rlorgat@math.harvard.edu}

\date{}

\begin{abstract}
For finite-dimensional evaluation modules with thin
$q$-characters, the Gautam--Toledano Laredo transfer from quantum
loop algebras to Yangians gives the Clebsch--Gordan input needed on
the evaluation-prefundamental Drinfeld--Kohno compact core. The
tensor product $V(a)\otimes L^-_i(b)$ has an associated graded
indexed by the $q$-character monomials of $V(a)$; if
$\chi_q(V)=\sum_m c_m m$, the corresponding prefundamental factor
appears with multiplicity $c_m$. Thin modules give the
coefficient-one form. For non-thin fundamental modules, already
present in classical types, the correct statement is the
multiset-valued filtration, not a multiplicity-free direct sum.
The unconditional output is the evaluation-generated compact-core
input to MC3; shifted-envelope generation and passage to the full
completed Drinfeld--Kohno category are separate completion
statements.
\end{abstract}

\subjclass[2020]{17B37, 17B69, 18M25, 81R10, 81R50}

\keywords{Drinfeld--Kohno category, Yangian, quantum affine algebra,
$q$-character, $\ell$-weight, Clebsch--Gordan decomposition,
Maurer--Cartan equation, chiral algebra, Chari--Moura}

\maketitle

% ================================================================
% 1. INTRODUCTION
% ================================================================
\section{Introduction}
\label{sec:intro}

\subsection{The minuscule hypothesis}

The classical Drinfeld--Kohno theorem \cite{Kohno87, Drinfeld90}
identifies the monodromy of the Knizhnik--Zamolodchikov
connection on conformal blocks for $\widehat{\fg}_k$ with the
$R$-matrix action of the quantum group
$U_q(\fg)$ at $q = e^{\pi i / (k + h^\vee)}$.
Kazhdan and Lusztig extended this to a braided tensor equivalence
between the Kazhdan--Lusztig category at level $k$ and the
category of type-$1$ modules over the quantum group
\cite{KL93a, KL93b, KL94a, KL94b}. Their proof, and the subsequent
refinement of Etingof and Kazhdan in terms of quantization of Lie
bialgebras \cite{EK96, EK98}, leans at several points on the
availability of a \emph{minuscule} representation: a fundamental
representation whose nonzero weight spaces are all one-dimensional
and form a single Weyl group orbit.

The minuscule hypothesis makes the block combinatorics transparent:
ordinary weight spaces are one-dimensional and are indexed by a
single Weyl orbit. Type~$A$ has this property for every fundamental
representation. Outside type~$A$ it is exceptional. In $E_8$ it is
absent: the adjoint representation is the smallest nontrivial
fundamental representation, its nonzero root spaces are
one-dimensional, and its zero-weight space has dimension
$\mathrm{rk}(E_8)=8$. A proof depending on minuscule fundamentals
therefore cannot be uniform in Lie type.

\subsection{Multiplicity-free $\ell$-weights}

\(q\)-characters refine ordinary characters by diagonalizing the
commuting Drinfeld loop-Cartan generators. The monomials of
\(\chi_q(L(\Psi))\) record the corresponding \(\ell\)-weights. For
thin modules, this refinement separates ordinary weight
multiplicities into coefficient-one monomials. For non-thin modules,
the coefficients of the \(q\)-character must be retained.

\begin{theorem}[Thin fundamental cases; Frenkel--Mukhin, Nakajima, Chari--Moura]
\label{thm:chari-moura-intro}
For the fundamental modules whose \(q\)-characters are thin,
\[
  \chi_q(V_{\omega_i}(a))=\sum_{s=1}^d m_s
\]
with pairwise distinct monomials and coefficient one. Nakajima's
quiver-variety computation~\cite{Nakajima04}, the
Frenkel--Mukhin algorithm~\cite{FrenkelMukhin01}, and the
Chari--Moura formulas~\cite{ChariMoura06} supply the source
computations in their respective ranges. In non-thin fundamental
cases the expansion has the form \(\chi_q(V)=\sum_m c_m m\) with
some \(c_m>1\).
\end{theorem}

The categorical input is therefore thinness where it holds, and a
multiset-valued \(q\)-character expansion where it does not.

\subsection{Statement of the main theorem}

Let $\Yh$ denote the Yangian of a simple Lie algebra $\fg$ with
formal parameter $\hbar$. For a fundamental weight $\omega_i$ and
spectral parameter $a \in \mathbb{C}$, let $V_{\omega_i}(a)$ be the
corresponding fundamental evaluation module, and let
$L^-_i(b) \in \cO^{\mathrm{sh}}$ denote the negative prefundamental
simple in Hern\'andez' shifted category, of highest $\ell$-weight
indexed by the Drinfeld polynomial of the identity representation
with spectral shift $b$ (see Definition~\ref{def:prefundamental}).
For a monomial $m$ in the $q$-character, write
$L^-_i(b;m)$ for the simple object whose highest $\ell$-weight is
obtained from that of $L^-_i(b)$ by multiplication by~$m$.

\begin{theorem}[Main theorem]
\label{thm:main}
For every finite-dimensional simple Lie algebra $\fg$, every
fundamental weight $\omega_i$, and all generic spectral parameters
$a, b \in \mathbb{C}$, the tensor product
$V_{\omega_i}(a)\otimes L^-_i(b)$ has a finite filtration in
$\cO^{\mathrm{sh}}$ with associated graded
\begin{equation}\label{eq:cg-main}
\gr\bigl(V_{\omega_i}(a) \otimes L^-_i(b)\bigr)
 \;\cong\;
 \bigoplus_{m \in \operatorname{Supp}\chi_q(V_{\omega_i}(a))}
 L^-_i(b;m)^{\oplus c_m},
 \qquad
 \chi_q(V_{\omega_i}(a))=\sum_m c_m m .
\end{equation}
The sum is taken with the $q$-character coefficients. Thin
fundamental modules give the coefficient-one form; non-thin
fundamental modules retain the displayed multiplicities. If the
relevant $\Ext^1$ groups between all listed simple factors,
including repeated factors, vanish on the chosen generic locus,
the filtration splits and (\ref{eq:cg-main}) is a
Clebsch--Gordan direct sum. Thus the evaluation-prefundamental
compact core of $\DK_\fg$ has a finite filtered
Clebsch--Gordan input in every simple Lie type. Completion beyond
that core is a separate theorem or hypothesis.
\end{theorem}

The restriction to generic spectral parameters excludes the finite
union of values of $a-b$ at which extra Hern\'andez $A$-factor
extensions occur. The displayed formula is the associated-graded
statement; a direct sum is exactly the additional Ext-vanishing
condition in
Theorem~\ref{thm:main}.

Theorem~\ref{thm:main} replaces the minuscule hypothesis of
Kazhdan--Lusztig and of earlier type-$A$ treatments of the
categorical Clebsch--Gordan filtration (see
\cite[\S 11]{ChariPressley94} and the Francis--Gaitsgory
pro-nilpotent completion of factorization categories
\cite{FG12}). The minuscule hypothesis is a computational shortcut:
it forces classical weight multiplicities to be trivial. The
$q$-character is the correct invariant for the shifted category.
Where it is thin, the graded factors occur once; where it is not,
the coefficients record the repeated prefundamental factors.

\subsection{Compact-core MC3 input}

The chiral bar complex $\barB^{\mathrm{ch}}(\cA)$ of a vertex
algebra $\cA$ admits a Maurer--Cartan element
$\Theta_\cA \in \MC(\fg^{\mathrm{mod}}_\cA)$ whose degree-$2$
projection is the classical $r$-matrix and whose higher-degree
projections form the shadow tower; see
\cite[Vol.\ I]{LorgatMKD}. The third master claim (MC3) in
modular Koszul duality asserts that the degree-$3$
component of $\Theta_\cA$, when restricted to the
evaluation-generated core, determines all higher-degree
components through thick generation and pro-nilpotent completion
of the ambient factorization category.
Theorem~\ref{thm:main} removes the Lie-type obstruction on the
evaluation-prefundamental core: for every simple $\fg$, the
categorical Clebsch--Gordan step needed by the Yangian target is a
finite filtered statement whose associated graded is controlled by
the $q$-character coefficients. The claim that the degree-$3$
shadow projection determines all higher projections belongs to the
compact-core MC3 theorem together with its completion hypotheses,
not to the representation-theoretic input alone.

\subsection{Minuscule weights and $\ell$-weight separation}

Kazhdan and Lusztig~\cite{KL93a} treat the equivalence with
quantum group modules type by type, with the minuscule weights of
the classical types doing the combinatorial work. Etingof and
Kazhdan~\cite{EK96, EK98} replace this by a universal
quantization functor on Lie bialgebras; their proof is
type-uniform, but the passage to the Yangian or quantum affine
setting reintroduces the minuscule hypothesis through the
prefundamental Clebsch--Gordan decomposition of
Chari--Pressley~\cite{ChariPressley94}. Hern\'andez and Jimbo
\cite{HernandezJimbo12} introduced the shifted category
$\cO^{\mathrm{sh}}$ and proved a block criterion which is the
combinatorial engine for the filtration. Francis and Gaitsgory
\cite{FG12} developed the pro-nilpotent completion formalism that
controls the separate passage from a compact-core comparison to the
full shifted category.

Theorem~\ref{thm:chari-moura-intro} is the replacement for the
minuscule hypothesis in the proof of the filtered categorical
Clebsch--Gordan statement. In the language of the Lorgat
monograph~\cite{LorgatMKD}, it supplies the categorical CG input
of MC3 for every simple Lie type, including $E_8$, in which no
minuscule representation exists.

% ================================================================
% 2. PRELIMINARIES
% ================================================================
\section{Preliminaries}\label{sec:prelim}

\subsection{Lie bialgebras and Yangians}

Let $\fg$ be a finite-dimensional simple Lie algebra over
$\mathbb{C}$ with Cartan subalgebra $\fh$, root system $\Phi$,
simple roots $\alpha_1, \ldots, \alpha_\ell$, fundamental weights
$\omega_1, \ldots, \omega_\ell$, Weyl group $W$, dual Coxeter
number $h^\vee$, and invariant symmetric form $(\cdot,\cdot)$
normalized so that long roots have squared length~$2$.
The polynomial loop algebra $\fg[t]$ carries a canonical rational
Lie bialgebra structure whose cobracket
$\delta: \fg[t] \to \fg[t] \otimes \fg[t]$ is determined by the
level-stripped rational $r$-matrix
\begin{equation}\label{eq:r-matrix-level}
r_{\mathrm{rat}}(u) = \frac{\Omg}{u},
\qquad
\Omg = \sum_a I_a \otimes I^a,
\end{equation}
where $\{I_a\}$ is a basis of $\fg$ and $\{I^a\}$ its dual with
respect to $(\cdot, \cdot)$. The affine collision calculation uses
the level-dependent trace-form residue
$r^{\mathrm{coll}}_k(z)=k\Omega_{\mathrm{tr}}/z$; the Yangian
and KZ rational structure use (\ref{eq:r-matrix-level}) with the
usual coupling normalization. The two normalizations must not be
identified. At affine level $k=0$ the trace-form double-pole
collision residue vanishes, while the simple-pole Lie bracket
channel and the level-stripped rational Yangian structure remain
nonzero; see~\cite{LorgatMKD}.

\begin{remark}[Level normalization]
\label{rem:level-norm}
The symbol $k$ belongs to the affine Kac--Moody collision residue.
The Yangian rational $r$-matrix is level-stripped. Thus $k=0$ is
not an abelianization of $\fg[t]$ and does not annihilate the
Yangian Maurer--Cartan structure.
\end{remark}

Drinfeld~\cite{Drinfeld85, Drinfeld88} quantized $\fg[t]$ to the
Yangian $\Yh = Y_\hbar(\fg)$: a Hopf algebra deforming
$U(\fg[t])$ along the Lie bialgebra structure
(\ref{eq:r-matrix-level}). The Yangian is graded by
conformal weight (the eigenvalues of a derivation lifting the
loop rotation $t \mapsto t + c$) and admits an evaluation
homomorphism $\mathrm{ev}_a: \Yh \to U(\fg)$ for each
$a \in \mathbb{C}$ in type~$A$, and more subtly in other types.

\subsection{The Drinfeld--Kohno category}

\begin{definition}[Drinfeld--Kohno categories]
\label{def:dk}
The finite Drinfeld--Kohno category $\DK^{\mathrm{fd}}_\fg$ is the
braided monoidal category of finite-dimensional $\Yh$-modules,
with braiding induced by the rational $R$-matrix attached to
(\ref{eq:r-matrix-level}). The shifted category
$\cO^{\mathrm{sh}}$ of Hern\'andez and Jimbo
~\cite{HernandezJimbo12} is larger: it contains the
infinite-dimensional prefundamental modules and has locally finite
Cartan-loop action with weights in a shifted dominant cone. In this
paper $\DK_\fg^{\mathrm{ec}}$ denotes the evaluation-prefundamental
compact core inside the completed shifted category, generated by
fundamental evaluation modules and negative prefundamental simples.
\end{definition}

\begin{definition}[Fundamental and prefundamental objects]
\label{def:prefundamental}
For each fundamental weight $\omega_i$ and spectral parameter
$a \in \mathbb{C}$, the \emph{fundamental evaluation module}
$V_{\omega_i}(a)$ is the finite-dimensional simple $\Yh$-module
whose Drinfeld polynomial at node $i$ is $1 - a z$ and is
trivial at all other nodes. The \emph{negative prefundamental
module} $L^-_i(b)$ is the simple infinite-dimensional object in
$\cO^{\mathrm{sh}}$ whose highest $\ell$-weight is the
``minus-delta at node $i$ with shift $b$'' configuration of
\cite{HernandezJimbo12}.
\end{definition}

\subsection{The chromatic filtration}

The category $\cO^{\mathrm{sh}}$ carries a decreasing
\emph{chromatic} or \emph{conformal-weight} filtration
$\cF^{\geq N} \subset \cO^{\mathrm{sh}}$, where $\cF^{\geq N}$
consists of objects whose generalized conformal-weight generalized
eigenvalues are bounded below by $N$. The associated graded
pieces $\cF^{\geq N} / \cF^{\geq N+1}$ are finite-dimensional
strata in the sense that, at each fixed highest $\ell$-weight,
the space of conformal-weight-$N$ vectors is finite-dimensional.
The compact evaluation-prefundamental core consists of the objects
supported on finitely many such strata and built from fundamental
evaluation modules and negative prefundamental simples by finite
extensions, shifts, cones, and retracts. Identifying this core with
the whole compact subcategory of the completed Drinfeld--Kohno
category is a separate completion theorem or hypothesis; it is not
part of the Clebsch--Gordan statement proved here.

\subsection{$\ell$-weights and the multiplicity-free property}

Let $\Uq$ denote the quantum affine algebra at the corresponding
parameter~$q$. Its commutative Cartan-type subalgebra is generated
by the loop elements $\phi_{i,\pm}(z)$ of Drinfeld's second
realization, and any finite-dimensional simple $\Uq$-module
decomposes as a direct sum of joint generalized eigenspaces for
these loop generators.

\begin{definition}[$\ell$-weights]
\label{def:l-weight}
An \emph{$\ell$-weight} of $\Uq$ is a pair
$(\boldsymbol{\psi}^+(z), \boldsymbol{\psi}^-(z))$ of
$\ell$-tuples of formal power series in $z^{\pm 1}$ subject to
the compatibility and rationality conditions of Frenkel and
Reshetikhin~\cite{FR99}. The \emph{$q$-character} of a
finite-dimensional simple module $V$ is
\[
\chi_q(V)
 \;=\;
 \sum_{m} \dim V_m \cdot m,
\]
where the sum ranges over $\ell$-weights $m$ and $V_m$ is the
corresponding joint generalized eigenspace. A module $V$ is
\emph{multiplicity-free in $\ell$-weights} if $\dim V_m \leq 1$
for every $\ell$-weight $m$.
\end{definition}

The central structural data we import from the literature is the
finite $q$-character with its coefficients.

\begin{theorem}[$q$-character coefficient data]
\label{thm:chari-moura}
For every finite-dimensional simple Lie algebra $\fg$, every
fundamental weight $\omega_i$, and every spectral parameter $a$,
the fundamental evaluation module $V_{\omega_i}(a)$ of $\Uq$ has a
finite $q$-character
\[
\chi_q(V_{\omega_i}(a))
 \;=\;
 \sum_m c_m m,
 \qquad
 c_m=\dim V_m .
\]
In the thin cases all coefficients are $1$. In the non-thin cases
the coefficients are retained and become multiplicities of the
prefundamental graded factors. Frenkel--Reshetikhin
\cite{FR99} define the $q$-character, Nakajima's
quiver-variety computation~\cite{Nakajima04} supplies the
simply-laced source calculation, and Chari--Moura
\cite{ChariMoura06} give explicit formulas for the classical
families in their range.
\end{theorem}

\subsection{Transfer between Yangians and quantum affine algebras}

The bridge between the Yangian $\Yh$ and the quantum affine algebra
$\Uq$ is the theorem of Gautam and Toledano Laredo:

\begin{theorem}[Gautam--Toledano Laredo, \cite{GTL17}]
\label{thm:gtl}
After choosing the logarithmic branch and avoiding the singular
hyperplanes of the parameters, there is a meromorphic tensor
equivalence
\[
\Phi_{\mathrm{GTL}} \;:\;
\Rep^{\mathrm{fd}}(\Uq) \;\xrightarrow{\;\sim\;}\;
\Rep^{\mathrm{fd}}(\Yh)
\]
between the corresponding finite-dimensional module categories.
It sends simples to simples and identifies the Drinfeld rational
fraction data on both sides. In particular,
$\Phi_{\mathrm{GTL}}$ preserves the $q$-character coefficients.
\end{theorem}

Consequently, Theorem~\ref{thm:chari-moura} transfers to the
Yangian as coefficient data: the fundamental evaluation module
$V_{\omega_i}(a)$ of $\Yh$ has the same finite multiset of
$\ell$-weights. Thin modules remain coefficient-one; non-thin
modules retain their multiplicities. The Hern\'andez--Jimbo block
criterion used below is combinatorial in the Drinfeld rational
fractions and therefore transfers on the same branch.

% ================================================================
% 3. STATEMENT
% ================================================================
\section{Statement of the main theorem}\label{sec:statement}

\begin{theorem}[Filtered categorical Clebsch--Gordan, all types]
\label{thm:main-formal}
Let $\fg$ be any finite-dimensional simple Lie algebra over
$\mathbb{C}$, let $\omega_i$ be any fundamental weight, and let
$a, b \in \mathbb{C}$ be spectral parameters lying in the
finite-filtration locus of $\cO^{\mathrm{sh}}$. Let
\[
\chi_q(V_{\omega_i}(a))=\sum_m c_m m,
\]
and write $L^-_i(b;m)$ for the negative prefundamental simple
obtained from $L^-_i(b)$ by multiplying its highest $\ell$-weight
by~$m$.
Then:
\begin{enumerate}[label=\textup{(\arabic*)}]
\item The tensor product $V_{\omega_i}(a)\otimes L^-_i(b)$ has a
finite filtration in $\cO^{\mathrm{sh}}$ whose associated graded is
\begin{equation}\label{eq:main-cg}
\gr\bigl(V_{\omega_i}(a) \otimes L^-_i(b)\bigr)
 \;\cong\;
 \bigoplus_{m\in \operatorname{Supp}\chi_q(V_{\omega_i}(a))}
 L^-_i(b;m)^{\oplus c_m}.
\end{equation}
\item The number of graded factors counted with multiplicity is
$\sum_m c_m=\dim V_{\omega_i}$. Thin modules are precisely the
coefficient-one case.
\item If
$\Ext^1_{\cO^{\mathrm{sh}}}(S,S')=0$ for all simple factors
$S,S'$ occurring in the multiset
\[
\{\,L^-_i(b;m) \text{ repeated } c_m \text{ times}\,\},
\]
then the filtration splits and (\ref{eq:main-cg}) is a
direct-sum Clebsch--Gordan decomposition.
\item The object $V_{\omega_i}(a)\otimes L^-_i(b)$ belongs to the
evaluation-prefundamental compact core generated by fundamental
evaluation modules and negative prefundamental simples. Identifying
that core with the full compact subcategory of a completed
Drinfeld--Kohno category is a separate completion theorem or
hypothesis.
\end{enumerate}
\end{theorem}

\begin{remark}[Scope of ``generic'']
\label{rem:generic-scope}
The generic locus excluded in the statement is the union of
complex hypersurfaces on which additional Hern\'andez $A$-factor
extensions appear. On the finite-filtration locus the associated
graded formula is the theorem. On the semisimple sublocus cut out
by the Ext-vanishing condition in
Theorem~\ref{thm:main-formal}(3), the filtration splits.
\end{remark}

% ================================================================
% 4. PROOF
% ================================================================
\section{Proof of the main theorem}\label{sec:proof}

The argument separates the character identity, the Yangian transfer,
the highest-$\ell$-weight filtration, and the Ext-vanishing
splitting criterion.

\subsection{Character identity}

\begin{lemma}[$q$-character identity, all types]
\label{lem:character-identity}
For every simple Lie algebra $\fg$, every fundamental weight
$\omega_i$, and all spectral parameters $a, b$:
\begin{equation}\label{eq:char-identity}
\bigl[V_{\omega_i}(a) \otimes L^-_i(b)\bigr]
 \;=\;
 \sum_{m\in \operatorname{Supp}\chi_q(V_{\omega_i}(a))}
 c_m\,[L^-_i(b;m)],
\end{equation}
in the completed Grothendieck group of the shifted category, where
$\chi_q(V_{\omega_i}(a))=\sum_m c_m m$.
\end{lemma}

\begin{proof}
The $q$-character records the joint generalized eigenspaces of the
Drinfeld Cartan currents. Tensoring a highest-$\ell$-weight
prefundamental object by a finite-dimensional module shifts the
highest $\ell$-weight by each monomial of the finite
$q$-character, with coefficient equal to the corresponding
$\ell$-weight dimension. This is the Hern\'andez--Jimbo
highest-$\ell$-weight calculus in
$\cO^{\mathrm{sh}}$~\cite{HernandezJimbo12}. The resulting
identity in the completed Grothendieck group is
(\ref{eq:char-identity}).
\end{proof}

\subsection{$q$-character coefficients for Yangian fundamentals}

\begin{lemma}[$q$-character coefficients for Yangian fundamentals]
\label{lem:qcharacter-yangian}
For every simple $\fg$ and every fundamental weight $\omega_i$,
the Yangian fundamental evaluation module $V_{\omega_i}(a)$
(viewed as a $\Yh$-module in $\cO^{\mathrm{sh}}$) has the same
$q$-character coefficient multiset as its quantum-loop counterpart
on the Gautam--Toledano Laredo branch.
\end{lemma}

\begin{proof}
For the quantum affine algebra $\Uq$, Theorem~\ref{thm:chari-moura}
gives the finite coefficient data. For the Yangian $\Yh$, apply
the Gautam--Toledano Laredo equivalence
$\Phi_{\mathrm{GTL}}$ of Theorem~\ref{thm:gtl}. This equivalence
is exact, sends simples to simples, and identifies the Drinfeld
rational fractions governing the commuting Cartan currents.
Therefore the joint generalized eigenspace dimensions, equivalently
the $q$-character coefficients, are preserved~\cite{GTL17}.
\end{proof}

\subsection{Highest-$\ell$-weight filtration}

\begin{lemma}[Filtration by prefundamental factors]
\label{lem:prefundamental-filtration}
Let $\chi_q(V_{\omega_i}(a))=\sum_m c_m m$, and let
$b \in \mathbb{C}$ be a spectral parameter in the finite-filtration
locus. Then
$V_{\omega_i}(a)\otimes L^-_i(b)$ admits a finite filtration whose
graded factors are the simples
$L^-_i(b;m)$ with multiplicity $c_m$.
\end{lemma}

\begin{proof}
The shifted category $\cO^{\mathrm{sh}}$ is controlled by
highest $\ell$-weights and the $A_{j,c}$ dominance order of
Hern\'andez--Jimbo~\cite{HernandezJimbo12}. Tensoring a
prefundamental simple by a finite-dimensional module produces a
finite object in this highest-$\ell$-weight order. Its possible
highest $\ell$-weights are obtained by multiplying the highest
$\ell$-weight of $L^-_i(b)$ by the monomials in
$\chi_q(V_{\omega_i}(a))$. Lemma~\ref{lem:qcharacter-yangian}
gives these monomials and their coefficients on the Yangian side.
The standard highest-$\ell$-weight filtration in
$\cO^{\mathrm{sh}}$ therefore has precisely the listed simple
graded factors, each repeated according to its $q$-character
coefficient.
\end{proof}

\subsection{Splitting criterion}

\begin{lemma}[Ext-vanishing splits the filtration]
\label{lem:ext-splitting}
Let $M$ be an object of $\cO^{\mathrm{sh}}$ with a finite
filtration whose graded factors are $S_1,\ldots,S_d$. If
$\Ext^1_{\cO^{\mathrm{sh}}}(S_j,S_k)=0$ for all $j,k$, then
$M\cong\bigoplus_j S_j$.
\end{lemma}

\begin{proof}
Split the filtration inductively. At each stage the extension class
lies in an $\Ext^1$ group between one graded factor and the direct
sum of the preceding factors. The stated vanishing kills that
class, so the short exact sequence splits. Iteration gives the
direct sum.
\end{proof}

\begin{proof}[Proof of Theorem~\ref{thm:main-formal}]
Set $M = V_{\omega_i}(a) \otimes L^-_i(b)$. Then $M$ is an object
of $\cO^{\mathrm{sh}}$: the tensor factor $V_{\omega_i}(a)$ is
finite-dimensional, $L^-_i(b)$ is in $\cO^{\mathrm{sh}}$ by
\cite{HernandezJimbo12}, and tensoring with a
finite-dimensional module preserves $\cO^{\mathrm{sh}}$.
Lemma~\ref{lem:prefundamental-filtration} gives the finite
filtration and identifies its graded factors. The number of factors
counted with multiplicity is $\sum_m c_m=\dim V_{\omega_i}$ by
Lemma~\ref{lem:qcharacter-yangian}.
Lemma~\ref{lem:ext-splitting} gives the split Clebsch--Gordan
decomposition on the Ext-vanishing sublocus. Finally, the
evaluation-prefundamental compact core is by definition the thick
idempotent-complete subcategory generated by fundamental evaluation
modules and negative prefundamental simples. Since thick
subcategories are closed under tensoring with generators, finite
extensions, cones, and retracts, the tensor product and its listed
graded factors lie in that core.
\end{proof}

% ================================================================
% 5. MC3 FOR ALL SIMPLE TYPES
% ================================================================
\section{Compact-core MC3 input}\label{sec:mc3}

Let $\cA$ be a vertex algebra and let
$\barB^{\mathrm{ch}}(\cA)$ denote its chiral bar complex on
ordered configuration space. Recall that the bar complex is
built on the augmentation ideal $\bar\cA = \ker(\varepsilon)$, not
on $\cA$ itself; we write
\[
\barB^{\mathrm{ch}}(\cA)
 \;=\; T^c(s^{-1} \bar\cA) \otimes
 \cO(\overline{C}_\bullet(X)),
\]
where $\cO(\overline{C}_n(X))$ is the sheaf of rational functions
on the Fulton--MacPherson compactification with logarithmic
poles along the diagonal divisor. The bar complex carries a
Maurer--Cartan element
\[
\Theta_\cA \in \MC(\fg^{\mathrm{mod}}_\cA)
\]
whose degree-$r$ projection $\pi_{r,0}(\Theta_\cA)$ is the
shadow datum at depth $r$; see~\cite{LorgatMKD} for the detailed
construction.

\begin{corollary}[Compact-core MC3 input for all simple types]
\label{cor:mc3-all-types}
Let $\fg$ be any finite-dimensional simple Lie algebra and let
$\cA_\fg = V_k(\fg)$ be the affine vertex algebra at level $k$
away from the critical level $k = -h^\vee$. On the thick closure
of fundamental evaluation modules and negative prefundamental
simples in $\DK_\fg$, the finite filtered Clebsch--Gordan input
required by the compact-core MC3 argument is available in every
simple Lie type. The assertion that the degree-$3$ projection
$\pi_{3,0}(\Theta_{\cA_\fg})$ determines all higher projections
then requires the separate compact-core MC3 theorem and its
completion hypotheses.
\end{corollary}

\begin{proof}
The degree-$2$ projection $\pi_{2,0}(\Theta_{\cA_\fg})$ is the
Yangian rational $r$-matrix in the level-stripped normalization
(\ref{eq:r-matrix-level}); the affine vertex algebra has collision
residue $k\Omega_{\mathrm{tr}}/z$. The two normalizations agree
only after the level convention is fixed. The degree-$3$ projection
satisfies the classical Yang--Baxter equation together with its
first Massey-product obstruction, but the passage from degree~$3$
to degree $r \geq 4$ is the separate pro-nilpotent completion
problem.

By Theorem~\ref{thm:main-formal}, for every simple Lie type
including $E_8$, the evaluation-prefundamental compact core has the
required filtered Clebsch--Gordan input; on the Ext-vanishing
semisimple locus this input is a direct-sum decomposition. The
lift from Grothendieck-group Baxter TQ relations to derived exact
triangles in the completed antidominant shifted envelope is the
completion statement of~\cite[\S 8]{LorgatMKD}. The
Francis--Gaitsgory pro-nilpotent formalism~\cite{FG12} then
controls the separate passage from the compact core to the
full ind-completion. The representation-theoretic obstruction is
the availability of this finite filtered $q$-character input, and
Theorem~\ref{thm:main-formal} supplies it for all simple types.
\end{proof}

\begin{remark}[The role of the level]
\label{rem:level-mc3}
The affine level $k$ enters the collision residue
$r^{\mathrm{coll}}_k(z)=k\Omega_{\mathrm{tr}}/z$ and the
Sugawara scalar, not the level-stripped Yangian rational
$r$-matrix (\ref{eq:r-matrix-level}). At $k=0$ the trace-form
double-pole residue vanishes, but the simple-pole Lie bracket
channel and the Yangian MC structure remain. The critical level
$k=-h^\vee$ is instead a singular Feigin--Frenkel centre regime
and is excluded from Corollary~\ref{cor:mc3-all-types}.
\end{remark}

% ================================================================
% 6. EXAMPLES
% ================================================================
\section{Examples by simple type}\label{sec:examples}

The input in Theorem~\ref{thm:main-formal} is the
$q$-character coefficient multiset. Minuscule fundamentals give
coefficient-one formulas. Non-minuscule fundamentals require the
actual coefficients of $\chi_q(V_{\omega_i}(a))$.

\subsection{Type $A_n$}

For $\fg=\fsl_{n+1}$, every fundamental representation is
minuscule. The $q$-character support is coefficient-one, and on
the Ext-vanishing locus the associated graded formula becomes the
usual split Clebsch--Gordan decomposition indexed by the weights
of $V_{\omega_i}$.

\subsection{Classical types $B_n,C_n,D_n$}

The minuscule fundamentals are the spin representation in type
$B_n$, the defining representation in type $C_n$, and the vector
and two half-spin representations in type $D_n$. These are the
coefficient-one cases. The remaining fundamentals are treated by
the Chari--Moura $q$-character formulas~\cite{ChariMoura06}. The
associated graded in Theorem~\ref{thm:main-formal} uses the
coefficients $c_m$ from those formulas; it is not an all-types
multiplicity-free statement.

\subsection{Exceptional types}

In $E_6$ the two minuscule fundamentals are coefficient-one. In
$E_7$ the unique minuscule fundamental is coefficient-one. The
other exceptional fundamentals, and every fundamental of $E_8$,
$F_4$, and $G_2$, are governed by their $q$-character coefficient
data. For $E_8$, the adjoint fundamental has an ordinary zero
weight space of dimension equal to the rank, so ordinary weights
cannot supply a multiplicity-free Clebsch--Gordan formula. The
filtered statement retains the coefficients of
$\chi_q(V_{\omega_i}(a))$ and then applies the prefundamental
filtration. A split formula requires the Ext-vanishing hypothesis
of Theorem~\ref{thm:main-formal}.

\begin{remark}[$q$-character input across simple types]
\label{rem:cg-table}
The following table records the input used in
Theorem~\ref{thm:main-formal}.
\begin{center}
\begin{tabular}{cccc}
\toprule
Type & Minuscule count & Key input & Output \\
\midrule
$A_n$ & $n$ & minuscule coefficient-one support & filtered core \\
$B_n$ & $1$ & Chari--Moura coefficients & filtered core \\
$C_n$ & $1$ & Chari--Moura coefficients & filtered core \\
$D_n$ & $3$ & Chari--Moura coefficients & filtered core \\
$E_6$ & $2$ & Nakajima coefficients & filtered core \\
$E_7$ & $1$ & Nakajima coefficients & filtered core \\
$E_8$ & $0$ & Nakajima coefficients & filtered core \\
$F_4$ & $0$ & $q$-character coefficients & filtered core \\
$G_2$ & $0$ & $q$-character coefficients & filtered core \\
\bottomrule
\end{tabular}
\end{center}
The minuscule hypothesis (all ordinary weight multiplicities
trivial) fails for the exceptional types $E_8$, $F_4$, $G_2$ and
for the intermediate fundamentals of the classical types. The
$q$-character coefficient multiset is the uniform replacement.
\end{remark}

% ================================================================
% 7. CONCLUSION
% ================================================================
\section{Conclusion}
\label{sec:conclusion}

The absence of minuscule representations in $E_8$ is a recurring
technical obstruction in the categorical parts of the quantum
geometric Langlands programme. Traditional proofs construct
equivalences of braided monoidal categories type by type, and
$E_8$ is either omitted or handled by a folding or restriction
procedure that reduces the representation theory to a setting with
minuscule weights. The coefficient-aware $q$-character filtration
supplies the replacement input on the Drinfeld--Kohno
evaluation-prefundamental core.

For the quantum geometric Langlands correspondence of
Gaitsgory and collaborators, the categorical CG step appears
in the proof that the spectral side (quasi-coherent sheaves on
the moduli of local systems) and the automorphic side (D-modules
on $\mathrm{Bun}_G$) are related by a functor of factorization
categories thickly generated by fundamental evaluation objects.
Corollary~\ref{cor:mc3-all-types} gives the Lie-type input for
that construction on the evaluation-prefundamental compact core,
including $E_8$. The existence of the full functor on completed
ambient categories remains the separate completion problem.

In modular Koszul duality, this supplies the
type-uniform compact-core Clebsch--Gordan input for MC3. The
pro-nilpotent passage from compact-core filtrations to the full
ambient ind-category is governed by the completion formalism of
Francis--Gaitsgory and by the module-category framework of
\cite{LorgatMKD}.

\begin{remark}[Limitations]
\label{rem:limitations}
Three caveats accompany Theorem~\ref{thm:main-formal}. First, the
generic spectral parameter hypothesis is genuine: on the
Hern\'andez $A$-linked hypersurfaces extra extensions may appear.
Second, the theorem proves an associated-graded filtered statement;
the direct-sum Clebsch--Gordan formula requires the stated
Ext-vanishing. Third, the theorem covers the
evaluation-prefundamental compact core; passage to completed
categories requires the Francis--Gaitsgory formalism and the
module-category framework of~\cite{LorgatMKD}.
\end{remark}

\begin{remark}[Comparison with prior literature]
\label{rem:comparison}
Drinfeld's ``Quasi-Hopf algebras''~\cite{Drinfeld90} introduced the
Drinfeld associator and the universal monodromy. Kazhdan--Lusztig
\cite{KL93a, KL93b, KL94a, KL94b} proved the categorical
equivalence type by type in positive-level affine categories.
Etingof--Kazhdan~\cite{EK96, EK98} gave a type-uniform
quantization. Hern\'andez and Jimbo~\cite{HernandezJimbo12}
introduced the shifted category $\cO^{\mathrm{sh}}$ and the
block criterion used here. Chari--Pressley~\cite{ChariPressley94}
proved the type-$A$ categorical CG. Francis--Gaitsgory
\cite{FG12} supplied the factorization pro-nilpotent formalism.
The $q$-character coefficient filtration is the uniform replacement
for the minuscule hypothesis, including $E_8$.
\end{remark}

% ================================================================
% BIBLIOGRAPHY
% ================================================================
\begin{thebibliography}{99}

\bibitem{ChariMoura06}
V.~Chari and A.~A.~Moura,
Characters and blocks for finite-dimensional representations of
quantum affine algebras,
\textit{Int.\ Math.\ Res.\ Not.}\ \textbf{2005} (2005), no.~5,
257--298.

\bibitem{ChariPressley94}
V.~Chari and A.~Pressley,
\textit{A Guide to Quantum Groups},
Cambridge University Press, 1994.

\bibitem{Drinfeld85}
V.~G.~Drinfeld,
Hopf algebras and the quantum Yang--Baxter equation,
\textit{Soviet Math.\ Dokl.}\ \textbf{32} (1985), 254--258.

\bibitem{Drinfeld88}
V.~G.~Drinfeld,
A new realization of Yangians and quantized affine algebras,
\textit{Soviet Math.\ Dokl.}\ \textbf{36} (1988), 212--216.

\bibitem{Drinfeld90}
V.~G.~Drinfeld,
Quasi-Hopf algebras,
\textit{Leningrad Math.\ J.}\ \textbf{1} (1990), no.~6, 1419--1457.

\bibitem{EK96}
P.~Etingof and D.~Kazhdan,
Quantization of Lie bialgebras I,
\textit{Selecta Math.\ (N.S.)}\ \textbf{2} (1996), 1--41.

\bibitem{EK98}
P.~Etingof and D.~Kazhdan,
Quantization of Lie bialgebras II,
\textit{Selecta Math.\ (N.S.)}\ \textbf{4} (1998), 213--231.

\bibitem{FeiginFrenkel92}
B.~Feigin and E.~Frenkel,
Affine Kac--Moody algebras at the critical level and
Gelfand--Dikii algebras,
\textit{Internat.\ J.\ Modern Phys.\ A}\ \textbf{7} (1992),
197--215.

\bibitem{FG12}
J.~Francis and D.~Gaitsgory,
Chiral Koszul duality,
\textit{Selecta Math.\ (N.S.)}\ \textbf{18} (2012), 27--87.

\bibitem{FR99}
E.~Frenkel and N.~Reshetikhin,
The $q$-characters of representations of quantum affine
algebras and deformations of $\cW$-algebras,
\textit{Contemp.\ Math.}\ \textbf{248} (1999), 163--205.

\bibitem{FrenkelMukhin01}
E.~Frenkel and E.~Mukhin,
Combinatorics of $q$-characters of finite-dimensional
representations of quantum affine algebras,
\textit{Comm.\ Math.\ Phys.}\ \textbf{216} (2001), 23--57.

\bibitem{GTL17}
S.~Gautam and V.~Toledano Laredo,
Meromorphic tensor equivalence for Yangians and quantum loop
algebras,
\textit{Publ.\ Math.\ IH\'ES}\ \textbf{125} (2017), 267--337.

\bibitem{HernandezJimbo12}
D.~Hern\'andez and M.~Jimbo,
Asymptotic representations and Drinfeld rational fractions,
\textit{Compos.\ Math.}\ \textbf{148} (2012), 1593--1623.

\bibitem{KL93a}
D.~Kazhdan and G.~Lusztig,
Tensor structures arising from affine Lie algebras I,
\textit{J.\ Amer.\ Math.\ Soc.}\ \textbf{6} (1993), 905--947.

\bibitem{KL93b}
D.~Kazhdan and G.~Lusztig,
Tensor structures arising from affine Lie algebras II,
\textit{J.\ Amer.\ Math.\ Soc.}\ \textbf{6} (1993), 949--1011.

\bibitem{KL94a}
D.~Kazhdan and G.~Lusztig,
Tensor structures arising from affine Lie algebras III,
\textit{J.\ Amer.\ Math.\ Soc.}\ \textbf{7} (1994), 335--381.

\bibitem{KL94b}
D.~Kazhdan and G.~Lusztig,
Tensor structures arising from affine Lie algebras IV,
\textit{J.\ Amer.\ Math.\ Soc.}\ \textbf{7} (1994), 383--453.

\bibitem{Kohno87}
T.~Kohno,
Monodromy representations of braid groups and Yang--Baxter
equations,
\textit{Ann.\ Inst.\ Fourier}\ \textbf{37} (1987), 139--160.

\bibitem{LorgatMKD}
R.~Lorgat,
\textit{Modular Koszul Duality, Volume I},
monograph, 2026.

\bibitem{Nakajima04}
H.~Nakajima,
Quiver varieties and $t$-analogs of $q$-characters of quantum
affine algebras,
\textit{Ann.\ of Math.\ (2)}\ \textbf{160} (2004), 1057--1097.

\end{thebibliography}

\end{document}
